\section{Smart Contracts as Move Generators}
\label{sec:sec06}
\label{sec:contracts}
%==============================================================================

A financial instrument is a deterministic program that emits moves, not an opaque record in a database. A bond emits coupons and a redemption; an option a premium and a conditional payoff; a swap periodic net payments. The logic varies; the output is always one list of moves.

\subsection{Contract Definition}

\begin{definition}[Smart Contract]
A smart contract is a deterministic function from inputs, state, and observed conditions to a sequence of atomic moves:
\[
\textit{Contract} : (\textit{Input}, \textit{State}, \textit{Conditions}) \to [\textit{Move}].
\]
\end{definition}

\noindent Equal inputs, state, and conditions yield equal moves: no randomness, no hidden state, no dependency beyond the passed parameters. The contract is not stored apart from the ledger---its effects \emph{are} ledger entries. Product complexity is isolated in the contract; the ledger sees only moves. A new product type is a new contract; the move stream, balance computation, and reporting are untouched, and its moves are indistinguishable from any other product's at the ledger level.

A smart contract is the machine-readable, executable form of the legal confirmation. An ISDA derivatives confirmation already states the product, economics, lifecycle events, and settlement obligations in structured terms; the Common Domain Model (CDM) is the formal vocabulary of that structure. Writing the contract from the confirmation is transcription, not engineering: CDM enums map directly onto product types, lifecycle events, and settlement obligations. The contract adds executability and nothing to the information content.

The output type is a \texttt{Move} carrying a quantity \texttt{Qty}. A quantity is an exact integer in the unit's smallest denomination (minor units: cents, satoshis, basis-point-notional); never a decimal, never a float. Integer minor units make addition associative and every balance computation exact, so identical inputs produce bit-identical outputs across hardware and compilers, the precondition for reproducible time travel and cross-system verification. Rounding (round-half-to-even, IEEE~754) applies once, at instruction projection---where an exact internal amount becomes a settlement or payment message---never within balance or payment arithmetic.

The types make this exactness structural. A \texttt{Qty} is an integer count of minor units, and the contract is a total function emitting a list of moves---the categorical shape is a fold over a finite schedule, meaning the contract is fully determined by the moves it lists, with no hidden step.

\begin{lstlisting}[language=Haskell]
-- A quantity is an EXACT integer in the unit's smallest denomination
-- (minor units: cents, satoshis, basis-point-notional). Never a decimal float.
newtype Qty = Qty Integer deriving (Eq, Ord, Show)

-- Qty is the additive abelian group of Section 2 (the position algebra).
instance Semigroup Qty where Qty a <> Qty b = Qty (a + b)  -- associative, exact
instance Monoid    Qty where mempty = Qty 0                 -- identity (zero)
qneg :: Qty -> Qty                                          -- inverse
qneg (Qty n) = Qty (negate n)

newtype WalletId = WalletId String deriving (Eq, Ord, Show)
newtype UnitId   = UnitId   String deriving (Eq, Ord, Show)

-- A move: a positive Qty of one unit, one wallet to another (Section 2),
-- tagged with the emitting contract and metadata. Direction carries the sign,
-- so mQty is always > 0.
data Move = Move
  { mFrom   :: WalletId, mTo     :: WalletId
  , mUnit   :: UnitId,   mQty    :: Qty       -- exact minor units, > 0
  , mTime   :: Integer,  mSource :: String    -- timestamp; emitting contract id
  , mMeta   :: String                         -- event description / external refs
  } deriving (Eq, Show)

-- A smart contract is a DETERMINISTIC, TOTAL function emitting a move schedule.
-- The ledger sees only [Move]; product complexity stays inside the contract.
type Contract i s c = (i, s, c) -> [Move]
\end{lstlisting}

\noindent Integer addition is associative and carries no rounding error, so balance and payment arithmetic is exact and identical inputs give bit-identical outputs. Rounding enters \emph{once}, at instruction projection, where converting an exact internal amount to a settlement message can introduce a fraction (a rate, an FX cross); round-half-to-even applies there and nowhere else.

\begin{lstlisting}[language=Haskell]
-- The SOLE rounding site. Prelude `round` IS round-half-to-even (IEEE 754).
projectAmount :: Rational -> Qty
projectAmount r = Qty (round r)

-- ghci> Qty 250 <> Qty 12500            -- exact integer addition
-- Qty 12750
-- ghci> qneg (Qty 12500) <> Qty 12500   -- conservation: no rounding error
-- Qty 0
-- ghci> projectAmount (5 % 2)           -- 2.5 -> nearest even -> 2
-- Qty 2
-- ghci> projectAmount (7 % 2)           -- 3.5 -> nearest even -> 4
-- Qty 4
-- ghci> Move (WalletId "wB") (WalletId "wS") (UnitId "USD")
--            (Qty 250000) 0 "put_ref" "PUT_PREMIUM"   -- Move 1 below, typed
\end{lstlisting}

\begin{remark}[FORMALIS clearance]
The \texttt{Qty} group (\texttt{(<>)}, \texttt{mempty}, \texttt{qneg}) is reused verbatim from the FORMALIS-cleared reference \texttt{reference/Ledger.hs} and is already cleared. The \texttt{Move} record, \texttt{Contract} synonym, and \texttt{projectAmount} are new: FORMALIS verdict \textbf{CLEARED}---all total (\texttt{round} is total on \texttt{Rational}; \texttt{Move} is a plain product; \texttt{Contract} is a type synonym) and deterministic, with rounding confined to the single \texttt{projectAmount} boundary so no balance arithmetic can round.
\end{remark}

\subsection{Composition: Modular over Composite}

\begin{principle}[Modular contracts]
One smart contract per payout or obligation type, not one composite contract per instrument.
\end{principle}

\noindent An option payoff depends on strike, spot, and exercise terms; a CSA margin depends on aggregate mark-to-market, threshold, minimum transfer amount, and eligible-collateral currency. These are independent calculations on independent inputs. Separation keeps each a pure function testable in isolation; combination makes the input space the product of both, most of which is irrelevant. A margin contract written once attaches to any collateral wallet---option, swap, or structured note---eliminating duplication.

Modularity requires clear interfaces: the executor---the single component permitted to mutate the ledger (Section~\ref{sec:lifecycle})---orchestrates the contract invocations for a given lifecycle event. This orchestration is simpler than maintaining a composite script for every product--CSA combination.

\subsection{Equity, Dividend, and Corporate-Action Contracts}

A stock contract emits moves for trade execution, dividend payment, and corporate actions through one primitive---a set of moves in a transaction---with no separate machinery.

\noindent\textbf{Trade execution.} A purchase of $n$ shares at price $p$ is a transaction of two atomic moves: $n \cdot p$ cash from buyer to seller, and $n$ of the stock from seller to buyer.

\noindent\textbf{Dividend.} On the ex-dividend date the contract emits cash from an issuer virtual wallet to each holder wallet, proportional to shares held. The unit state records the event (date, amount per share) for idempotency.

\noindent\textbf{Bond accrued interest.} Accrued interest is not a separate unit. The bond price function $P_t(u)$ returns the \emph{dirty} price (including accrual), and a between-coupon trade settles as cash-for-bond at the dirty price. The next coupon date pays the full coupon and transitions the bond ex-coupon. Clean versus dirty is a reporting convention derived from unit state (last coupon date, accrual convention), not a separate position; both prices are reported by subtracting accrued interest computed from that state.

\noindent\textbf{Corporate actions.} A corporate action is a sequence of state transitions and moves anchored to its dates: announcement, record, ex/effective, payment. At the record date the unit state records the entitlement snapshot; at the ex/effective date the contract emits position-adjustment moves from a corporate-action virtual wallet---a 2-for-1 split doubles each entitled holding while the price reference halves. Each step is a transaction with full metadata (dates, ratios, entitlement rules), so time travel reconstructs both pre- and post-action views. The same multi-date mechanism handles mergers, spin-offs, and rights issues.

\begin{remark}[Unit state is a projection, not a mutable store]
Unit state above is materialised \emph{UnitStatus}, the read cache governed in Section~\ref{sec:lifecycle}. UnitStatus holds one value per unit, shared---read identically by every holder. Its value changes over time, but UnitStatus is not a separate source of truth: the immutable event log is. Every change is caused by a logged event; the stored value is overwritten in place only as a read cache. Replaying the events up to any point rebuilds the exact value that held then, so nothing is lost by overwriting and there is no other way to change the value. A value exists from registration onward.
\end{remark}

\subsection{European Put---Full Move Schedule}

\noindent Parameters: strike $K$, expiry $T$, underlying $S$, notional $N$, premium $P$, buyer $w_B$, seller $w_S$. The contract is completely defined by conditional moves at inception and expiry.

\noindent\textit{Move 1 (inception, $t=0$)} --- unconditional, on contract creation:

\begin{lstlisting}
Move(from: w_B, to: w_S, unit: USD, quantity: P * N,
     timestamp: t_0, source: "put_ref", metadata: "PUT_PREMIUM")
\end{lstlisting}

\noindent\textbf{Cash settlement.} \textit{Move 2 (expiry $t=T$)} executes only if $S_T < K$ (in-the-money); the seller pays the buyer the intrinsic value $(K-S_T)\cdot N$. If $S_T \ge K$, no move.

\begin{lstlisting}
Move(from: w_S, to: w_B, unit: USD, quantity: (K - S_T) * N,
     timestamp: T, source: "put_ref", metadata: "PUT_SETTLEMENT_CASH")
\end{lstlisting}

\noindent\textbf{Physical delivery.} If $S_T < K$, the buyer delivers $N$ of the underlying and receives $K\cdot N$ cash; the two moves are one atomic transaction, requiring the buyer to hold the underlying at expiry. If $S_T \ge K$, no moves.

\begin{lstlisting}
Move(from: w_B, to: w_S, unit: S,   quantity: N,     ...   metadata: "PUT_DELIVERY_ASSET")
Move(from: w_S, to: w_B, unit: USD, quantity: K * N, ...   metadata: "PUT_DELIVERY_CASH")
\end{lstlisting}

\noindent\textbf{Lot-size constraint.} Equity delivery requires whole multiples of the board lot $L$ (reference data in the Unit Store). For lot $L$,
\[
N_{\text{deliver}} = \left\lfloor \tfrac{N}{L} \right\rfloor \cdot L, \qquad N_{\text{residual}} = N - N_{\text{deliver}}.
\]
Whole lots settle as a DvP pair at strike; the residual settles in cash at intrinsic value. All moves are one atomic transaction:

\begin{lstlisting}
Transaction (settlement: txIntroduce = Nothing, carrying the three moves below):
    Move(from: w_B, to: w_S, unit: S,   quantity: N_deliver)            -- whole lots
    Move(from: w_S, to: w_B, unit: USD, quantity: K * N_deliver)        -- strike cash
    Move(from: w_S, to: w_B, unit: USD, quantity: (K - S_T) * N_residual) -- cash residual
\end{lstlisting}

\begin{proposition}[Lot-split conservation]
The underlying nets to zero ($-N_{\text{deliver}} + N_{\text{deliver}} = 0$); cash nets to zero, the seller's outflow $K\cdot N_{\text{deliver}} + (K-S_T)\cdot N_{\text{residual}}$ equalling the buyer's inflow. The economic value transferred is the put payoff $(K-S_T)\cdot N$, independent of the lot split.
\end{proposition}

\noindent The lot-adjustment is a pure function of $(K, N, S_T, L)$, tested by the property suite (Appendix~\ref{app:properties}). The settlement projection (Section~\ref{sec:settlement}) yields two instructions: a securities delivery (\texttt{sese.023}) for $N_{\text{deliver}}$ shares and a cash payment for $K\cdot N_{\text{deliver}} + (K-S_T)\cdot N_{\text{residual}}$. The pattern generalises to any physically settled instrument with lot constraints---convertible conversion, warrant exercise, equity-forward delivery---only the payoff formula differs.

The contract is fully defined by its move schedule: a finite list of conditional transfers parameterised by market observables. Pricing complexity (Black--Scholes, Monte Carlo, finite differences) is irrelevant to the ledger, which sees only the emitted moves.

\subsection{Interest Rate Swap---Full Move Schedule}

\noindent Parameters: notional $N$, fixed rate $r_{\text{fixed}}$, floating reference $r_{\text{float},k}$ (e.g.\ SOFR) at reset $k$, reset dates $\{t_1,\ldots,t_n\}$, day-count fraction $\delta_k$ (CDM date types and conventions fix the dates and $\delta_k$, Section~\ref{sec:cdm-dates}), fixed-payer $w_F$, floating-payer $w_L$. For each reset $t_k$:
\[
\text{Payment}_k = N \cdot (r_{\text{float},k} - r_{\text{fixed}}) \cdot \delta_k .
\]

\begin{lstlisting}
Move(from: w_F, to: w_L, unit: USD, quantity: Payment_k,
     timestamp: t_k, source: "irs_ref", metadata: "IRS_NET_PAYMENT")
\end{lstlisting}

\noindent Move quantities are positive (Section~\ref{sec:ledger}). The contract emits $|\text{Payment}_k|$: $w_F \to w_L$ when $\text{Payment}_k > 0$, $w_L \to w_F$ when $\text{Payment}_k < 0$. Direction encodes the sign; quantity is the magnitude. Each reset needs only the current floating rate and the fixed parameters---no historical path. The complete contract is the sequence $\{m_1,\ldots,m_n\}$, one move per reset, plus the unit state recording settled resets for idempotency.

\medskip
\noindent\textbf{Worked example (CDM lifecycle).} A 5-year USD IRS, notional \$10M, fixed 3\%, quarterly resets.

\begin{itemize}[nosep]
\item \emph{Inception} (CDM \texttt{ExecutionEvent}, intent \texttt{OPEN}). Unit state initialised: fixed rate, floating index, schedule, reset counter $=0$. No cash move (par swap, no upfront premium).
\item \emph{Reset at $t_1$} (CDM \texttt{ResetEvent}). SOFR observed at 3.5\%; unit state records \texttt{float\_rate\_$t_1$} $=3.5\%$, counter incremented.
\item \emph{Coupon at $t_1$} (CDM \texttt{TransferEvent}). $10\text{M}\times(3.5\%-3\%)\times0.25 = 12{,}500$. \texttt{Move(from: w\_F, to: w\_L, unit: USD, quantity: 12{,}500)}; period $t_1$ marked settled.
\item \emph{Maturity at $t_n$} (CDM \texttt{TerminationEvent}, intent \texttt{TERMINATE}). Final coupon as above; unit state transitions to TERMINATED. No further resets or payments.
\end{itemize}

\clearpage
%==============================================================================
